\documentclass[12pt]{article}

\usepackage{fontawesome}
\usepackage{hyperref}
\usepackage{xurl}
\usepackage[tagged, highstructure]{accessibility}
\usepackage{bookmark}
\usepackage{enumitem}

\hypersetup{
    colorlinks=false,
    pdfborder={0 0 0},
    pdftitle={Data Visualization Lab},
    pdfauthor={Adrianna Holden-Gouveia},
    pdfsubject={Introduction to Data - Lab Assignment},
    pdflang={en-US},
    pdfkeywords={data visualization; charts; graphs; accessibility; storytelling with data}
}

% Configure PDF bookmarks for navigation
\bookmarksetup{
    numbered,
}

% Configure list spacing for better accessibility
\setlist{nosep}



\title{Data Visualization}
\author{
        Adrianna Holden-Gouveia \\
        Website: \href{https://aholdengouveia.name}{aholdengouveia.name}\\
        %Email: \href{mailto:admin@aholdengouveia.name}{admin@aholdengouveia.name} \\
%        \faLinkedin{: aholdengouveia} \\
        \faGithub {: aholdengouveia} \\
        %\faTwitter {: aholdengouveia} \\
        }
\date{\vspace{-5ex}}

%\date{\today}


\begin{document}    

\maketitle

%\begin{abstract}
%This is a template for Lab work
%\end{abstract}
%\tableofcontents


\section*{Accessibility Notice}
This document is also available in HTML format at:

\href{https://aholdengouveia.name/IntroData/labs/datavisualizations.html}
{\textbf{\nolinkurl{https://aholdengouveia.name/IntroData/labs/datavisualizations.html}}}

The HTML version provides enhanced accessibility features including keyboard navigation, screen reader support, responsive design, dark mode support, and high contrast options.

\section*{Objectives:}
\begin{itemize}
    \item The goal of this lab assignment is for students to distinguish between good and bad data visualizations using real-world examples. Through practical examples and observations, students will gain an understanding of when to use which data visualizations and how best to implement them.
\end{itemize}
\section*{Complete the following problems}

References, a video, a PowerPoint and some notes are available at

\href{https://www.aholdengouveia.name/IntroData/datavis.html}
{\textbf{\nolinkurl{https://www.aholdengouveia.name/IntroData/datavis.html}}}

Examples of good Data Visualizations
\begin{itemize}
    \item Information is Beautiful \href{https://informationisbeautiful.net/}{\textbf{\nolinkurl{https://informationisbeautiful.net/}}}
    \item Storytelling with data \href{https://www.storytellingwithdata.com/blog}{\textbf{\nolinkurl{https://www.storytellingwithdata.com/blog}}}
\end{itemize}    

Examples of bad Data Visualizations
\begin{itemize}
    \item WTF Visualizations: Visualizations that make no sense \href{https://viz.wtf/}{\textbf{\nolinkurl{https://viz.wtf/}}}
    \item r/dataisugly \href{https://www.reddit.com/r/dataisugly/}{\textbf{\nolinkurl{https://www.reddit.com/r/dataisugly/}}}
\end{itemize}



\subsection*{Evaluate visualizations}
Using the guidelines from this website \href{https://ixdf.org/literature/article/guidelines-for-good-visual-information-representations}{\textbf{\nolinkurl{https://ixdf.org/literature/article/guidelines-for-good-visual-information-representations}}} and the materials from this week, answer the following questions.  For questions 1 to 3, use the table and checklist (DataVizChecklist.docx, linked with this lab) to break down your answer and make sure it hits each of the guidelines.

    \begin{enumerate}
        \item Your favorite visualization
        \begin{enumerate}
            \item Include a screenshot of the visualization you picked and its URL.
            \item Why is it your favorite?
            \item Fill out the table and checklist for it.
        \end{enumerate}
        \item Your least favorite or most aggravating visualization
        \begin{enumerate}
            \item Include a screenshot of the visualization you picked and its URL.
            \item Why is it your least favorite?
            \item Fill out the table and checklist for it.
        \end{enumerate}
        \item The visualization you found most humorous. This can come from any of the example websites.
        \begin{enumerate}
            \item Include a screenshot of the visualization you picked and its URL.
            \item Why did you find it the most humorous?
            \item Fill out the table and checklist for it.
        \end{enumerate}
        \item Accessibility: from the good examples, pick one you like that you haven't already evaluated. We're going to look at the accessibility of it.
        \begin{enumerate}
            \item Include a screenshot of the visualization you picked and its URL.
            \item Name 3 ways you can make this visualization more accessible.
            \item Write an appropriate alt text for it.
            \item Create a narration of the visualization. Your narration can be written or a video. One easy way to make a video is to create a screen capture video using something like OBS, upload it to a site like YouTube, and make it unlisted so you can just share a link.
        \end{enumerate}
        \item Go to \href{https://www.economist.com/graphic-detail}{\textbf{\nolinkurl{https://www.economist.com/graphic-detail}}}, do you think in general, this link belongs in the good or bad category? You don't need to evaluate every visualization, just look through the first page or two, and explain in a paragraph which one you think this link belongs in, why, and include screenshots of the visualizations you picked that support your opinion.
    \end{enumerate}


    \subsection*{Create your own Visualizations}
    \begin{enumerate}
        \item Going back to the Halloween candy data, pick 1 thing that you think would be better illustrated using a visualization.  Create your visualization using whatever software you like, this could be an online chart maker, or a spreadsheet or anything. 
        \item Explain why you picked that 1 thing, and explain how you decided which visualization to use and why.
        \item Go back to the data you collected for the Data Collection lab in Week 2, and pick 2 things that you think would be better illustrated using a visualization. Create your visualization using whatever software you like, this could be an online chart maker, or a spreadsheet or anything. 
        \item Explain why you picked those 2 things from your data, and explain how you decided which visualization to use and why.
    \end{enumerate}

    \section*{Deliverables}
    \begin{sloppypar}
    \begin{enumerate}
        \item A document with answers to each of the above questions evaluating other people's visualizations. Make sure to include a screenshot and the URL of each visualization you picked, and fill out the linked table and checklist for each of the visualizations in questions 1 to 3.
        \item A document for your visualizations, the document should include screenshots (do not just upload multiple images to Brightspace) of the 3 visualizations that you created. Each screenshot should have the explanation of why below it.
    \end{enumerate}
    \end{sloppypar}
\end{document}